\begin{figure}[H]
\centering
\includegraphics[width=1.0\linewidth]{hokkaido/F7-1_wind_roses.pdf}
\caption[風配図（前半）]{代表格子点の風配図、\textbf{月の前半}（1--15 日）。
平均と楕円では見せられない「分布そのもの」を示す。
ローブが 2 つに割れている場面が、平均風がほぼゼロで分散だけが大きくなる時期にあたる。}
\label{fig:f7-hokkaido-h1}
\end{figure}

\begin{figure}[H]
\centering
\includegraphics[width=1.0\linewidth]{hokkaido/F7-2_wind_roses.pdf}
\caption[風配図（後半）]{同じ格子点の\textbf{月の後半}（16 日--月末）。
前半（図\ref{fig:f7-hokkaido-h1}）と並べると、\textbf{2 枚はよく似ている}。
\SI{30}{\hecto\pascal} の卓越風向は 1・4・7・10 月のいずれも前半と後半で同じである。
これは第 \ref{sec:validation} 章の結論と整合する。
月をさらに細かく割っても、新しい構造は出てこない。}
\label{fig:f7-hokkaido-h2}
\end{figure}